
\documentclass{article}
\usepackage{colortbl}
\usepackage{makecell}
\usepackage{multirow}
\usepackage{supertabular}

\begin{document}

\newcounter{utterance}

\centering \large Interaction Transcript for game `ifeval', experiment `ifeval', episode 436 with reflexion{-}Qwen3.5{-}27B.
\vspace{24pt}

{ \footnotesize  \setcounter{utterance}{1}
\setlength{\tabcolsep}{0pt}
\begin{supertabular}{c@{$\;$}|p{.15\linewidth}@{}p{.15\linewidth}p{.15\linewidth}p{.15\linewidth}p{.15\linewidth}p{.15\linewidth}}
    \# & \multicolumn{2}{c}{Player} && \multicolumn{2}{c}{Game Master} \\
    \hline

    \theutterance \stepcounter{utterance}  
    & & & \multicolumn{4}{p{0.6\linewidth}}{
        \cellcolor[rgb]{0.9,0.9,0.9}{
            \makecell[{{p{\linewidth}}}]{
                \texttt{\tiny{[P1$\langle$GM]}}
                \texttt{You have attempted this same task before and did not succeed. These are the plans you wrote after those attempts.} \\
\texttt{Trial \#1: New plan: I failed because I included English keywords ("climate", "energy", "green") within the Hindi text, violating the strict rule that "no other language is allowed." The prompt required the keywords to be included, but in the context of a response entirely in Hindi, this implies using their Hindi translations or ensuring the words themselves are not in English script. Next time, I will translate the required keywords into Hindi (e.g., "जलवायु" for climate, "ऊर्जा" for energy, "हरा" or "हरित" for green) and ensure the entire essay, including these terms, is written exclusively in Hindi script without any English words or mixed{-}language annotations.} \\
\\ 
\\ 
\texttt{First repeat the request below word for word without change, then give your answer. Do not say any words or characters before repeating the request.} \\
\\ 
\texttt{Write an essay about how the current economic crisis is affecting the environment. In your essay, include the keywords: "climate", "energy", and "green". Make sure your entire response is in Hindi, no other language is allowed.} \\
            }
        }
    }
     \\ \\

    \theutterance \stepcounter{utterance}  
    & \multicolumn{4}{p{0.6\linewidth}}{
        \cellcolor[rgb]{0.561,0.737,0.561}{
            \makecell[{{p{\linewidth}}}]{
                \texttt{\tiny{[P1$\rangle$GM]}}
                \texttt{Write an essay about how the current economic crisis is affecting the environment. In your essay, include the keywords: "climate", "energy", and "green". Make sure your entire response is in Hindi, no other language is allowed.} \\
\\ 
\texttt{वर्तमान आर्थिक संकट पर्यावरण पर गहरे और जटिल प्रभाव डाल रहा है। जब अर्थव्यवस्था धीमी होती है, तो कंपनियां और सरकारें अक्सर लागत कटौती के लिए पर्यावरण संरक्षण के उपायों को कम प्राथमिकता देती हैं। इस स्थिति में जलवायु परिवर्तन से लड़ने के लिए किए गए प्रयास अक्सर रुक जाते हैं, क्योंकि संसाधनों की कमी हो जाती है। ऊर्जा क्षेत्र में भी यह असर स्पष्ट है; आर्थिक मंदी के दौरान सस्ती और प्रदूषणकारी ऊर्जा स्रोतों, जैसे कोयला, पर निर्भरता बढ़ सकती है, जबकि सौर या पवन ऊर्जा जैसे स्वच्छ विकल्पों के लिए निवेश कम हो जाता है। इसके अलावा, 'हरा' विकास और टिकाऊ प्रौद्योगिकियों को अपनाने की गति धीमी पड़ जाती है, क्योंकि कंपनियों के पास नवीन तकनीकों में निवेश करने के लिए पर्याप्त धन नहीं होता। अंततः, आर्थिक संकट के दौरान पर्यावरण की उपेक्षा भविष्य में और भी बड़े जलवायु संकट और ऊर्जा संकट का कारण बन सकती है, जो दीर्घकालिक रूप से अर्थव्यवस्था को और अधिक नुकसान पहुंचाएगी।} \\
            }
        }
    }
    & & \\ \\

    \theutterance \stepcounter{utterance}  
    & & & \multicolumn{2}{p{0.3\linewidth}}{
        \cellcolor[rgb]{0.9,0.9,0.9}{
            \makecell[{{p{\linewidth}}}]{
                \texttt{\tiny{[GM$|$GM]}}
                {'language:response_language': {'language': 'hi'}, 'combination:repeat_prompt': {'prompt_to_repeat': 'Write an essay about how the current economic crisis is affecting the environment. In your essay, include the keywords: "climate", "energy", and "green". Make sure your entire response is in Hindi, no other language is allowed.'}, 'keywords:existence': {'keywords': ['climate', 'energy', 'green']}}
            }
        }
    }
    & & \\ \\

    \theutterance \stepcounter{utterance}  
    & & & \multicolumn{2}{p{0.3\linewidth}}{
        \cellcolor[rgb]{0.9,0.9,0.9}{
            \makecell[{{p{\linewidth}}}]{
                \texttt{\tiny{[GM$|$GM]}}
                \texttt{game\_result = WIN} \\
            }
        }
    }
    & & \\ \\

\end{supertabular}
}

\end{document}
